\begin{frame}[fragile]{Stellingen en Bewijzen (Opmaak)}
    Er zijn 3 standaard `vormen' voor een omgeving:
    \begin{itemize}[label={\textbullet}]
        \item \texttt{plain}: schuingedrukte tekst, extra ruimte om de omgeving heen 
            (vaak gebruikt in stellingen, lemmas, proposities)
        \pause
        \item \texttt{definition} rechtgedrukte tekst, extra ruimte om de omgeving heen 
            (vaak gebruikt in definities, voorbeelden, opgaves)
        \pause
        \item \texttt{remark} Rechtgedrukte tekst, \emph{geen} extra ruimte om de omgeving heen 
            (vaak gebruikt in opmerkingen, samenvattnigen, conclusies)
    \end{itemize}
    \pause 
    Je maakt een omgeving in deze stijl door:
    \begin{minted}{tex}
        \theoremstyle{plain}
        \newtheorem{stelling}{Stelling}
    \end{minted}
\end{frame}

\begin{frame}[fragile]{Stellingen en Bewijzen (Opmaak)}
    Om dus een aantal omgevingen goed te zetten, zowel in nummering als in stijl:
    \begin{minted}[highlightlines={1,6,10}]{tex}
        \theoremstye{plain}
        \newtheorem{stelling}{Stelling}[section]
        \newtheorem{lemma}[stelling]{Lemma}
        \newtheorem{propositie}[stelling]{Propositie}
        | \pause |
        \theoremstyle{definition}
        \newtheorem{definitie}[stelling]{Definitie}
        \newtheorem{voorb}[stellin]{Voorbeeld}
        | \pause |
        \theoremstyle{remark}
        \newtheorem*{opm}{Opmerking}
    \end{minted}    
\end{frame}